\documentclass[conference]{IEEEtran}
\IEEEoverridecommandlockouts

\usepackage{cite}
\usepackage{amsmath,amssymb,amsfonts}
\usepackage{algorithmic}
\usepackage{graphicx}
\usepackage{textcomp}
\usepackage{xcolor}
\usepackage{booktabs}
\usepackage{multirow}
\usepackage{array}
\usepackage{url}
\usepackage{hyperref}
\usepackage{microtype}

% Prevent aggressive hyphenation on long technical terms
\hyphenation{pre-dic-tive main-te-nance im-bal-ance XG-Boost Light-GBM SMOTE LSTM}

\def\BibTeX{{\rm B\kern-.05em{\sc i\kern-.025em b}\kern-.08em
    T\kern-.1667em\lower.7ex\hbox{E}\kern-.125emX}}

\begin{document}

\title{Hybrid Machine Learning and Deep Learning for Health Status Classification and Remaining Useful Life Prediction in Manufacturing}

\author{
Achmad Zikran Maulida\textsuperscript{1},
Amir Hamzah\textsuperscript{2},
Aqsa Zamzami\textsuperscript{3}, 
Reynaldi Chandra Kirana\textsuperscript{4} \\

\textsuperscript{}Jurusan Teknik Informatika dan Komputer, 
Program Studi Teknik Informatika \\
Politeknik Negeri Jakarta \\
Jln. Prof. DR. G.A. Siwabessy, Kampus Universitas Indonesia, Depok 16425 \\

\textsuperscript{1}achmad.zikran.maulida.tik23@stu.pnj.ac.id \\
\textsuperscript{2}amir.hamzah.tik23@stu.pnj.ac.id \\
\textsuperscript{3}aqsa.zamzami.tik23@stu.pnj.ac.id \\
\textsuperscript{4}reynaldi.chandra.kirana.tik23@stu.pnj.ac.id
}
\maketitle

\begin{abstract}
Machine failure events constitute as few as 0.056\% of industrial
IoT sensor observations, creating an extreme class imbalance that,
when compounded by temporal data leakage from naive train-test
splitting, renders most standard predictive maintenance pipelines
unreliable under real deployment conditions. This paper presents
PRIME (Predictive Reliability and Intelligence Maintenance Engine),
an end-to-end pipeline that resolves both challenges through three
principled contributions: (1)~\textbf{Stratified Sequential Block
Sampling (SSBS)}, a temporally-aware undersampling strategy that
reduces HEALTHY-class dominance from 97.4\% to 87.1\% while
preserving chronological data integrity; (2)~\textbf{Temporal Label
Engineering} with a P90-based Sensor Confirmation Layer that
converts binary failure signals into three operationally actionable
health states (HEALTHY, WARNING, CRITICAL); and (3)~a
\textbf{cascaded XGBoost--LSTM inference architecture} that
activates the regression module conditionally upon anomaly
detection, reducing real-time inference overhead by
$\sim$70\% compared to always-on sequential execution. Evaluated
on a 100,000-record, 20-machine time-series dataset under strict
machine-based splitting with train-only SMOTE augmentation, the
pipeline achieves a WARNING-class F1-score of 0.9818 with
\textbf{zero fatal misclassifications} (CRITICAL$\rightarrow$HEALTHY),
and predicts Remaining Useful Life with an MAE of 0.7985~days,
with 98.04\% of predictions falling within one day of ground truth.
These results establish that principled temporal handling of
severely imbalanced sensor data enables production-grade predictive
maintenance without sacrificing computational efficiency.
\end{abstract}

\begin{IEEEkeywords}
Predictive Maintenance, XGBoost, LSTM, Temporal Label Engineering, SSBS, Remaining Useful Life, Time-Series Classification, SMOTE
\end{IEEEkeywords}

%=============================================================
\section{INTRODUCTION}
%=============================================================

The emergence of Industry 4.0 has enabled pervasive deployment of Internet-of-Things (IoT) sensors across manufacturing facilities, generating continuous high-frequency streams of machine operational data \cite{lee2014industryiot}. Traditional maintenance paradigms—corrective (fix on failure) and preventive (fixed-schedule)—are increasingly inadequate: the former incurs costly unplanned downtime, while the latter wastes resources on unnecessary interventions. Predictive Maintenance (PdM) leverages sensor data and machine learning to schedule maintenance precisely when needed, transforming reactive operations into data-driven, condition-based workflows \cite{zhang2019pdmsurvey}.

Despite this promise, real-world PdM deployments face two compounding technical challenges. First, failure events are exceedingly rare in well-maintained industrial machinery. In the dataset analyzed in this work, failure events represent only 56 records out of 100,000 total observations---a failure rate of 0.056\%. This constitutes an extreme class imbalance problem, where standard classifiers are overwhelmingly biased toward predicting the majority HEALTHY class \cite{chawla2002smote}. Resampling techniques such as SMOTE (Synthetic Minority Over-sampling Technique) \cite{chawla2002smote} are commonly applied to mitigate this, but they introduce a critical secondary risk: if applied before train-test splitting, synthetic samples generated from future data contaminate the training set, causing optimistic bias known as temporal data leakage \cite{kaufman2012leakage}.

Second, the temporal structure of time-series sensor data imposes constraints that render standard cross-validation and random splitting invalid. Conventional $k$-fold cross-validation violates temporal causality by including future observations in training folds, leading to inflated performance estimates that do not generalize to real-world deployment \cite{bergmeir2018tscv}. Machine-based splitting—assigning entire machine histories to distinct splits—addresses this by ensuring each model learns to generalize across unseen machines, reflecting realistic inference conditions.

Existing literature lacks a unified framework that simultaneously addresses: (1) extreme class imbalance without temporal leakage, (2) multi-class label engineering from binary failure signals, and (3) a cascaded ML+DL inference architecture that is computationally efficient in real-time industrial deployment. This paper fills this gap by presenting the PRIME pipeline, with the following contributions:

\begin{enumerate}
\item We propose \textbf{Stratified Sequential Block Sampling (SSBS)}, a temporally-aware undersampling strategy that selects chronological blocks per machine, preserving time-series causality while reducing HEALTHY dominance from 97.36\% to 87.09\%.
\item We introduce a \textbf{three-class Temporal Label Engineering} scheme with an empirically calibrated Sensor Confirmation Layer (P90 thresholds, minimum 2 sensors), transforming binary failure labels into actionable HEALTHY/WARNING/CRITICAL classifications.
\item We demonstrate that \textbf{machine-based train/validation/test splitting combined with train-only SMOTE} constitutes an enterprise-grade zero-leakage pipeline for time-series PdM modeling.
\item We propose a \textbf{cascaded XGBoost+LSTM deployment architecture} that achieves state-of-the-art classification (F1=0.9906) and RUL prediction (MAE=0.7985 days, Error$\leq$1d=98.04\%) while reducing inference overhead by $\sim$70\% through conditional LSTM activation.
\end{enumerate}

%=============================================================
\section{RELATED WORK}
%=============================================================

\subsection{Deep Learning for RUL Prediction}
Long Short-Term Memory (LSTM) networks have become the dominant architecture for Remaining Useful Life prediction from sequential sensor data. Hochreiter and Schmidhuber's foundational LSTM \cite{hochreiter1997lstm} addressed the vanishing gradient problem inherent in vanilla RNNs, enabling learning of long-range temporal dependencies critical for degradation modeling. Zheng et al. \cite{zheng2017lstm_rul} demonstrated LSTM superiority over feed-forward networks on the CMAPSS turbofan dataset, achieving significant RMSE reductions. More recently, hybrid CNN-LSTM architectures \cite{li2018cnnlstm} have been proposed to extract local feature representations (CNN) before temporal modeling (LSTM), outperforming pure LSTM on multi-sensor industrial datasets. Gated Recurrent Units (GRU), introduced by Cho et al. \cite{cho2014gru}, offer a computationally lighter alternative to LSTM; however, empirical studies on complex multi-sensor degradation signals consistently favor LSTM due to its superior capacity for modeling asymmetric temporal patterns across different operational zones.

\subsection{Tree-Based Methods for Fault Classification}
Ensemble tree methods remain the state-of-the-art for tabular and time-series classification in industrial fault detection. Chen and Guestrin \cite{chen2016xgboost} introduced XGBoost, a scalable gradient-boosted tree framework that consistently outperforms alternatives on structured data. Its regularization parameters ($\alpha$, $\lambda$) and gradient-based pruning make it robust to overfitting in high-dimensional feature spaces generated by rolling statistics and lag features. Random Forest \cite{breiman2001rf}, while often used as a strong baseline, tends to produce wider decision boundaries that inflate false alarm rates in imbalanced industrial datasets. Ke et al. \cite{ke2017lightgbm} proposed LightGBM, achieving faster training via leaf-wise growth and histogram-based splitting, though its aggressive early stopping behavior can cause premature convergence when minority class samples are sparse, as observed in our experiments.

\subsection{Imbalance Handling in Industrial Data}
Class imbalance is ubiquitous in PdM datasets, where failure events are by design rare \cite{li2020pdmimbalance}. SMOTE \cite{chawla2002smote} generates synthetic minority samples by interpolating between existing instances in feature space, effectively augmenting training data without data collection overhead. However, applying SMOTE to the full dataset before splitting introduces temporal leakage: synthetic samples may encode statistical properties of future (test) data, producing overly optimistic evaluation results. Correct practice mandates applying SMOTE exclusively within the training fold after splitting \cite{blagus2013smote_bias}. ADASYN and Borderline-SMOTE variants improve on standard SMOTE by focusing synthesis near decision boundaries, but their benefits are marginal when threshold calibration is applied post-hoc as a complementary strategy.

\subsection{Time-Series Splitting Strategies}
Standard $k$-fold cross-validation is invalid for time-series data because it samples future observations into training folds, violating temporal causality \cite{bergmeir2018tscv}. Walk-forward (expanding window) validation is the canonical solution for univariate time series \cite{hyndman2018forecasting}. For multi-entity industrial datasets, machine-based splitting—where distinct machines are assigned to train, validation, and test partitions—is more appropriate, as it evaluates generalization across unseen machines rather than across time for the same machine. This reflects realistic deployment conditions where models must predict on new, previously unseen equipment. Existing literature often neglects this distinction, applying global temporal splits that concentrate minority-class failure events into the test partition and starve training of critical degradation patterns.

\subsection{Research Gap}
While individual contributions exist for LSTM-based RUL prediction, XGBoost fault classification, and SMOTE imbalance handling, no prior work proposes a unified pipeline that: (1) applies SSBS for temporally-safe undersampling before splitting, (2) enforces machine-based splits to prevent cross-machine leakage, (3) applies SMOTE exclusively on the training set post-split, (4) employs a cascaded ML+DL architecture for efficient deployment. The temporal leakage risk from early SMOTE application is frequently overlooked, and the computational efficiency of conditional LSTM activation has not been systematically studied in the PdM context.

%=============================================================
\section{DATASET AND PREPROCESSING}
%=============================================================

\subsection{Dataset Description}

The experimental dataset (\texttt{sensor\_readings.csv}) comprises continuous hourly IoT sensor readings collected from 20 identical industrial machines over 208 days (1 July 2025 to 25 January 2026). Each machine contributes exactly 5,000 timestamped records, yielding 100,000 total observations with 11 columns: one timestamp, one machine identifier, eight sensor measurements, and one binary failure label. The dataset exhibits extreme class imbalance: only 56 records are labeled \texttt{failure=1} (0.056\%), with an average of 2.80 failure events per machine. Eight sensor features span thermal, mechanical, electrical, and environmental domains: \texttt{temperature} ($^\circ$C), \texttt{vibration} (mm/s), \texttt{pressure} (bar), \texttt{rpm}, \texttt{power\_consumption} (kW), \texttt{noise\_level} (dB), \texttt{humidity} (\%), and \texttt{operating\_hours} (cumulative). Temporal integrity was verified: zero duplicate timestamps and zero inter-record gaps exceeding one hour. A supplementary \texttt{maintenance\_logs.csv} (500 records, 8 columns) provides technician notes used for NLP-derived feature enrichment.

\begin{table}[htbp]
\caption{TABLE I: Dataset Statistics}
\centering
\begin{tabular}{ll}
\hline
\textbf{Property} & \textbf{Value} \\
\hline
Total Records & 100,000 \\
Number of Machines & 20 \\
Observation Period & 208 days \\
Records per Machine & 5,000 \\
Sensor Features & 8 \\
Metadata Columns & 3 (timestamp, machine\_id, failure) \\
Failure Events & 56 (0.056\%) \\
Healthy Events & 99,944 (99.944\%) \\
Avg. Failures / Machine & 2.80 \\
Temporal Resolution & 1 hour \\
\hline
\end{tabular}
\label{tab:dataset_stats}
\end{table}

\subsection{Temporal Label Engineering}

The binary \texttt{failure} label is insufficient for operational decision-making, as it provides no advance warning of impending equipment degradation. We transform it into a three-class ordinal label: \textbf{HEALTHY} (normal operation), \textbf{WARNING} (degradation onset, actionable maintenance window), and \textbf{CRITICAL} (imminent failure, emergency intervention required).

Label assignment follows a temporal backward-labeling algorithm: for each \texttt{failure=1} event at timestamp $T_f$, all records in the look-back window $[T_f - W_{WARNING}, T_f)$ are labeled WARNING and records in $[T_f - W_{CRITICAL}, T_f)$ are relabeled CRITICAL (overriding WARNING), where $W_{WARNING} = 48$ hours and $W_{CRITICAL} = 24$ hours. These window parameters were derived empirically from Failure Autopsy analysis (Phase 2 EDA), examining sensor trajectories 72 hours prior to each failure event across machines M-01, M-05, and M-09. Analysis revealed consistent degradation signal onset at $T-48$h and dramatic escalation at $T-24$h, validating the chosen windows. This constitutes a revision from the initial blueprint hypothesis of $W_{WARNING} = 72$h.

A \textbf{Sensor Confirmation Layer} further refines the WARNING label to reduce false positives: a WARNING label is only retained if at least 2 of the 6 high-priority sensors (temperature, vibration, pressure, rpm, power\_consumption, noise\_level) exceed their respective P90 thresholds computed from the HEALTHY distribution. This confirmation step downgraded 49 WARNING rows to HEALTHY (1.82\% of WARNING candidates), while all CRITICAL rows were retained. The resulting label distribution is shown in Table~\ref{tab:label_dist} and P90 thresholds in Table~\ref{tab:p90}.

The health label formula for record $i$ of machine $m$ is:
\begin{equation}
y_i = \begin{cases}
2 & \text{if } T_f - W_{CRITICAL} \leq t_i < T_f \\
1 & \text{if } T_f - W_{WARNING} \leq t_i < T_f - W_{CRITICAL} \\
  & \quad \text{and } \sum_{s=1}^{6} \mathbf{1}[x_{i,s} > \theta_s^{P90}] \geq 2 \\
0 & \text{otherwise}
\end{cases}
\label{eq:label}
\end{equation}
where $\theta_s^{P90}$ is the 90th percentile threshold for sensor $s$ computed over HEALTHY rows, and $\mathbf{1}[\cdot]$ is the indicator function.

\begin{table}[htbp]
\caption{TABLE II: Final Label Distribution After Temporal Engineering}
\centering
\begin{tabular}{lrr}
\hline
\textbf{Class} & \textbf{Count} & \textbf{Percentage} \\
\hline
HEALTHY (0) & 97,364 & 97.364\% \\
WARNING (1) & 1,247 & 1.247\% \\
CRITICAL (2) & 1,389 & 1.389\% \\
\hline
\textbf{Total} & \textbf{100,000} & \textbf{100.000\%} \\
\hline
\end{tabular}
\label{tab:label_dist}
\end{table}

\begin{table}[htbp]
\caption{TABLE III: P90 Sensor Thresholds — Confirmation Layer}
\centering
\begin{tabular}{ll}
\hline
\textbf{Sensor} & \textbf{P90 Threshold (HEALTHY)} \\
\hline
temperature & 76.30 $^\circ$C \\
vibration & 0.59 mm/s \\
pressure & 103.80 bar \\
rpm & 2,540.00 \\
power\_consumption & 82.10 kW \\
noise\_level & 74.30 dB \\
\hline
\end{tabular}
\label{tab:p90}
\end{table}

%%====================================================================
%% FIGURE PLACEHOLDER B: Sensor Timeline Menuju Failure
%% Sumber gambar: notebooks/fase_2_eda/02_eda_forensik.ipynb
%% Cari cell yang memplot time-series sensor (temperature, vibration, dll)
%% di sekitar periode failure untuk mesin M-01/M-05/M-09.
%% Export gambar sebagai: figures/sensor_timeline_failure.png
%% Cara di Overleaf: Upload ke folder "figures/", uncomment blok di bawah.
%%====================================================================
% \begin{figure}[htbp]
%   \centering
%   \includegraphics[width=\columnwidth]{figures/sensor_timeline_failure.png}
%   \caption{Sensor trajectory for representative machine M-05 during the
%            72~hours preceding a failure event. Consistent degradation
%            onset is visible at $T{-}48$h (WARNING boundary) with rapid
%            escalation at $T{-}24$h (CRITICAL boundary), validating
%            the empirically derived window parameters.}
%   \label{fig:sensor_timeline}
% \end{figure}
%%====================================================================

\subsection{Sensor Informativeness Analysis}

Cohen's $d$ effect size was computed for each sensor to quantify the separation between HEALTHY and failure-zone (WARNING+CRITICAL) distributions:
\begin{equation}
d = \frac{\mu_{failure} - \mu_{healthy}}{\sqrt{\frac{(n_1-1)\sigma_1^2 + (n_2-1)\sigma_2^2}{n_1+n_2-2}}}
\label{eq:cohend}
\end{equation}

All six primary sensors exhibit large effect sizes ($d > 2.0$), confirming their high discriminative power for label assignment and feature prioritization. Operating hours and humidity exhibit medium effect sizes and were retained but deprioritized in feature importance analysis. Results are shown in Table~\ref{tab:cohend}.

\begin{table}[htbp]
\caption{TABLE IV: Cohen's D Effect Size per Sensor}
\centering
\begin{tabular}{llll}
\hline
\textbf{Sensor} & \textbf{Cohen's d} & \textbf{Category} & \textbf{Priority} \\
\hline
vibration & 3.37 & Large & High \\
pressure & 3.06 & Large & High \\
rpm & 2.91 & Large & High \\
noise\_level & 2.89 & Large & High \\
temperature & 2.71 & Large & High \\
power\_consumption & 2.61 & Large & High \\
operating\_hours & 0.29 & Medium & Low \\
humidity & 0.21 & Medium & Low \\
\hline
\end{tabular}
\label{tab:cohend}
\end{table}

%=============================================================
\section{FEATURE ENGINEERING}
%=============================================================

Raw sensor values alone provide limited temporal context for predictive modeling. We enrich the feature space through five systematic engineering groups, yielding a final feature matrix of 69 columns (after excluding label and identifier columns) as detailed in Table~\ref{tab:features}.

\textbf{Rolling Statistics} capture local temporal trends using two window sizes (24h and 48h) applied to the six high-priority sensors, computing mean, standard deviation, and maximum for each combination: $6 \times 2 \times 3 = 36$ features. These windows align with the empirically derived degradation onset periods and are the top-ranked features by XGBoost importance (e.g., \texttt{temperature\_roll\_max\_24h} ranks first).

\textbf{Lag Features} encode recent historical sensor readings at offsets of 1, 6, and 24 hours for the six high-priority sensors: $6 \times 3 = 18$ features. These capture autoregressive dependencies and enable the classifier to detect rate-of-change patterns indicative of accelerating degradation.

\textbf{Cross-Sensor Ratios} (4 features) encode nonlinear interaction effects between sensor pairs (e.g., vibration-to-pressure ratio), capturing multivariate degradation signatures not visible in univariate signals.

\textbf{Degradation Proxy} (1 feature): \texttt{hours\_since\_last\_maintenance}, computed by joining with \texttt{maintenance\_logs.csv}, quantifies cumulative wear since the last corrective or preventive intervention.

\textbf{NLP-Derived Features} (2 features): \texttt{damage\_category} (encoded categorical) and \texttt{severity\_score} (ordinal integer) are extracted from technician notes in \texttt{maintenance\_logs.csv} via keyword-based text mining and mapped to each machine's operational history.

All 69 numeric features are standardized using \texttt{StandardScaler} fit exclusively on training machines (M-01 to M-14) to prevent scaling-induced leakage. Negative vibration values ($<0$) introduced by sensor noise were clipped to zero prior to scaling.

\begin{table}[htbp]
\caption{TABLE V: Feature Engineering Summary (69 Input Features)}
\centering
\begin{tabular}{lll}
\hline
\textbf{Group} & \textbf{Count} & \textbf{Description} \\
\hline
Original Sensors & 8 & Raw sensor values \\
Rolling Statistics & 36 & 6 sensors $\times$ 2 windows $\times$ 3 stats \\
Lag Features & 18 & 6 sensors $\times$ lag \{1, 6, 24\}h \\
Cross-Sensor Ratios & 4 & Pairwise sensor interaction \\
Degradation Proxy & 1 & hours\_since\_last\_maintenance \\
NLP-Derived & 2 & damage\_category, severity\_score \\
\hline
\textbf{TOTAL} & \textbf{69} & After label/ID columns removed \\
\hline
\end{tabular}
\label{tab:features}
\end{table}

%=============================================================
\section{IMBALANCE HANDLING AND DATA SPLITTING}
%=============================================================

\subsection{Stratified Sequential Block Sampling (SSBS)}

Naive random undersampling of the HEALTHY class to reduce imbalance destroys temporal continuity: removing arbitrary rows breaks the sequential structure required by lag features and rolling statistics. Pure temporal undersampling (e.g., keeping only the first $N$ rows per machine) similarly fails, as it truncates the pre-failure degradation zone—the most informative period for learning.

We propose \textbf{Stratified Sequential Block Sampling (SSBS)}: for each of the 20 machines, two chronological blocks are selected:
\begin{itemize}
\item \textbf{Block A (first 500 rows)}: captures prime operating conditions when the machine is healthy and sensor signals are nominal.
\item \textbf{Block B (last 500 rows with 24h safety buffer)}: captures the pre-warning boundary period. A 24-hour safety buffer is applied before the first WARNING zone to avoid partial-label contamination at the block boundary.
\end{itemize}

All WARNING and CRITICAL rows are always retained in full (100\% inclusion) regardless of block boundaries. The target quota of 1,000 HEALTHY rows per machine (20 machines $\times$ 1,000 = 20,000) is partially achieved; machines with early failures have fewer available rows in Block A/B (e.g., M-06: 523, M-10: 534), resulting in an organic shortfall that reflects temporal integrity rather than a bug. The SSBS formula for HEALTHY selection per machine $m$ is:

\begin{equation}
S_m^{HEALTHY} = \text{Block}_A(m, 500) \cup \text{Block}_B(m, 500, \delta=24h)
\label{eq:ssbs}
\end{equation}

where $\delta$ is the safety buffer duration. The resulting dataset of 20,423 rows (Table~\ref{tab:ssbs}) preserves all WARNING and CRITICAL instances while substantially reducing HEALTHY dominance from 97.36\% to 87.09\%.

\begin{table}[htbp]
\caption{TABLE VI: SSBS Result Distribution}
\centering
\begin{tabular}{lrrl}
\hline
\textbf{Class} & \textbf{Before SSBS} & \textbf{After SSBS} & \textbf{After \%} \\
\hline
HEALTHY & 97,364 & 17,787 & 87.09\% \\
WARNING & 1,247 & 1,247 & 6.11\% \\
CRITICAL & 1,389 & 1,389 & 6.80\% \\
\hline
\textbf{TOTAL} & \textbf{100,000} & \textbf{20,423} & \textbf{100\%} \\
\hline
\end{tabular}
\label{tab:ssbs}
\end{table}

%%====================================================================
%% FIGURE PLACEHOLDER C: SSBS Block Selection Diagram
%% Sumber gambar: notebooks/fase_6_imbalance/06_imbalance_handling.ipynb
%% Cari cell yang memvisualisasikan distribusi class sebelum/sesudah SSBS,
%% atau diagram Block A dan Block B per mesin.
%% Export gambar sebagai: figures/ssbs_diagram.png
%% Cara di Overleaf: Upload ke folder "figures/", uncomment blok di bawah.
%%====================================================================
% \begin{figure}[htbp]
%   \centering
%   \includegraphics[width=\columnwidth]{figures/ssbs_diagram.png}
%   \caption{SSBS block selection strategy per machine. Block~A captures
%            prime operating conditions (rows 1--500) and Block~B
%            captures the pre-failure boundary (last 500 rows with a
%            24-hour safety buffer). All WARNING and CRITICAL rows are
%            retained unconditionally.}
%   \label{fig:ssbs_diagram}
% \end{figure}
%%====================================================================

\subsection{Machine-Based Train/Validation/Test Split}

A global temporal split (e.g., first 70\% of timestamps for training) is inappropriate for run-to-failure datasets: since WARNING and CRITICAL events cluster at the end of each machine's timeline, a temporal split concentrates minority-class samples into the test set, leaving training devoid of degradation patterns. We adopt a \textbf{machine-based split} that assigns complete machine histories to each partition:

\begin{itemize}
\item \textbf{Train:} Machines M-01 through M-14 (14 machines)
\item \textbf{Validation:} Machines M-15, M-16, M-17 (3 machines)
\item \textbf{Test:} Machines M-18, M-19, M-20 (3 machines)
\end{itemize}

This ensures each partition contains complete run-to-failure cycles, the model learns to generalize across unseen machines (realistic deployment), and zero temporal overlap exists between splits. The \texttt{StandardScaler} is fit exclusively on training machine data and applied (transform-only) to validation and test machines.

\begin{table}[htbp]
\caption{TABLE VII: Machine-Based Dataset Split}
\centering
\begin{tabular}{llrl}
\hline
\textbf{Split} & \textbf{Machines} & \textbf{Samples} & \textbf{Percentage} \\
\hline
Train & M-01 to M-14 & 14,419 & 70.60\% \\
Validation & M-15, M-16, M-17 & 3,288 & 16.10\% \\
Test & M-18, M-19, M-20 & 2,716 & 13.30\% \\
\hline
\textbf{Total} & 20 machines & \textbf{20,423} & \textbf{100\%} \\
\hline
\end{tabular}
\label{tab:split}
\end{table}

\subsection{SMOTE on Training Set Only}

Residual class imbalance in the training set (HEALTHY: 12,504 vs. WARNING: 901 vs. CRITICAL: 1,014) is addressed via SMOTE \cite{chawla2002smote}, applied \textit{exclusively} to the training partition after splitting. This strict ordering is critical: applying SMOTE before splitting would allow synthetic minority samples—generated by interpolating between existing minority instances—to encode statistical properties derived from validation or test machine data, constituting temporal and cross-machine data leakage \cite{kaufman2012leakage}.

SMOTE targets both WARNING and CRITICAL classes to 4,000 samples each, generating 3,099 synthetic WARNING and 2,986 synthetic CRITICAL instances. The HEALTHY class is not oversampled. Validation and test sets remain 100\% natural, containing no synthetic data, ensuring unbiased evaluation.

\begin{table}[htbp]
\caption{TABLE VIII: SMOTE Effect on Training Set}
\centering
\begin{tabular}{lrrr}
\hline
\textbf{Class} & \textbf{Before SMOTE} & \textbf{After SMOTE} & \textbf{Synthetic Added} \\
\hline
HEALTHY & 12,504 & 12,504 & 0 \\
WARNING & 901 & 4,000 & 3,099 \\
CRITICAL & 1,014 & 4,000 & 2,986 \\
\hline
\textbf{TOTAL} & \textbf{14,419} & \textbf{20,504} & \textbf{6,085} \\
\hline
\end{tabular}
\label{tab:smote}
\end{table}

%=============================================================
\section{PREDICTIVE MODELS}
%=============================================================

\subsection{Model 1: Health Status Classifier}

The Health Status Classifier operates as the first stage of the cascaded architecture, performing real-time multi-class inference on every incoming sensor tick. Three candidate algorithms were trained and evaluated on the SMOTE-augmented training set (20,504 samples, 69 features) and assessed on the natural validation and test sets.

\textbf{\textit{Random Forest} (Baseline):} An ensemble of 300 decision trees trained with \texttt{n\_estimators=300}, \texttt{max\_depth=20}, \texttt{min\_samples\_split=5}, \texttt{min\_samples\_leaf=2}, \texttt{max\_features='sqrt'}, and \texttt{class\_weight='balanced'}. This configuration establishes a strong baseline, achieving F1$_{macro}$=0.9914 on the test set. However, the validation F1$_{macro}$ of 0.9292 reveals a 0.0622 train-validation gap attributable to high-variance tree boundaries, and WARNING F1 on validation (0.811) indicates excessive false alarms under natural class distribution.

\textbf{\textit{XGBoost V2} (Selected):} A gradient-boosted tree ensemble trained with \texttt{n\_estimators=1000}, \texttt{max\_depth=4}, \texttt{learning\_rate=0.01}, \texttt{subsample=0.8}, \texttt{reg\_alpha=0.5}, \texttt{reg\_lambda=2.0}, converging at iteration 499 via early stopping. The initial V1 configuration (learning\_rate=0.05, early stopping patience=40) converged prematurely at iteration 40 (mlogloss=0.317 vs. 0.121 for V2), producing catastrophic WARNING F1=0.2186. V2 addresses this by tuning early stopping patience and reducing regularization aggressiveness. Additionally, a \textbf{WARNING probability threshold} of 0.60 (vs. default 0.50) is applied post-training, eliminating 953 false alarms (954$\rightarrow$1) while maintaining WARNING F1=0.9818 on validation.

\textbf{\textit{LightGBM} (Baseline):} A leaf-wise gradient boosting model with \texttt{learning\_rate=0.05} and \texttt{early\_stopping\_rounds=50}, which converged aggressively at iteration 27 due to logloss stagnation in early iterations under sparse minority samples. A WARNING threshold of 0.65 recovers validation F1$_{macro}$ from 0.6663 to 0.9845, making it competitive but inferior to XGBoost on the primary validation WARNING F1 metric.

\begin{table}[htbp]
\caption{TABLE IX: Model 1 Hyperparameter Summary}
\centering
\begin{tabular}{p{2.2cm}ccc}
\hline
\textbf{Parameter} & \textbf{Rand.~Forest} & \textbf{XGBoost V2} & \textbf{LightGBM} \\
\hline
n\_estimators   & 300          & 499$^{\dagger}$ & 27$^{\dagger}$ \\
max\_depth      & 20           & 4               & N/A \\
learning\_rate  & N/A          & 0.01            & 0.05 \\
subsample      & N/A          & 0.8             & N/A \\
reg\_alpha      & N/A          & 0.5             & N/A \\
reg\_lambda     & N/A          & 2.0             & N/A \\
class\_weight   & balanced     & N/A             & N/A \\
WARN threshold & 0.50         & \textbf{0.60}   & 0.65 \\
Model size     & 3.15 MB      & 1.68 MB         & 0.13 MB \\
\hline
\multicolumn{4}{l}{$^{\dagger}$ Converged via early stopping.}
\end{tabular}
\label{tab:hyper_clf}
\end{table}

\subsection{Model 2: RUL Predictor}

The RUL Predictor operates as the second stage, activated exclusively when Model 1 outputs WARNING or CRITICAL. This design is justified by four empirical observations:
\begin{enumerate}
\item \textbf{Irreducible uncertainty in HEALTHY zone:} Sensor signals during normal operation lack the discriminative patterns necessary for accurate long-horizon RUL estimation (MAE$_{HEALTHY}$=13.77 days vs. MAE$_{WARNING}$=0.10 days).
\item \textbf{Business relevance:} Factory operators require actionable short-horizon predictions (hours to days), not uncertain long-range estimates (90+ days).
\item \textbf{Architectural consistency:} Model 1 detection of WARNING/CRITICAL defines the decision boundary for Model 2 activation.
\item \textbf{Computational efficiency:} Conditional activation reduces LSTM inference frequency by $\sim$70\%, as HEALTHY constitutes $>$87\% of operational time.
\end{enumerate}

The RUL regression dataset comprises 1,915 training samples (WARNING: 901, CRITICAL: 1,014), 288 validation samples, and 433 test samples, filtered to WARNING+CRITICAL rows only. The RUL target variable \texttt{rul\_days} is engineered by computing forward time-to-next-failure for each row, with median values of $\sim$2 days for WARNING and $\sim$1 day for CRITICAL, consistent with the $W_{WARNING}=48$h and $W_{CRITICAL}=24$h windows.

\textbf{\textit{XGBoost Regressor}:} Trained with \texttt{n\_estimators=1000} (best\_iteration=497), \texttt{max\_depth=4}, \texttt{learning\_rate=0.01}, \texttt{subsample=0.8}, \texttt{colsample\_bytree=0.8}, \texttt{min\_child\_weight=3}, \texttt{gamma=0.1}, \texttt{reg\_alpha=0.3}, \texttt{reg\_lambda=1.5}, \texttt{objective=reg:squarederror}, \texttt{eval\_metric=mae}. The top feature by permutation importance is \texttt{power\_consumption\_roll\_std\_48h}, confirming domain validity.

\textbf{\textit{LSTM V2} (Selected):} A two-layer stacked LSTM network processing sequences of 24 consecutive timesteps ($\times$ 69 features). Training proceeded for 200 epochs (maximum), with the best checkpoint at epoch 184. Architecture:
\begin{itemize}
\item LSTM(64 units) + L2(0.001) + Dropout(0.3) + BatchNormalization
\item LSTM(32 units) + L2(0.001) + Dropout(0.3) + BatchNormalization
\item Dense(16, ReLU) $\rightarrow$ Dense(1, Linear)
\item Total parameters: 47,649
\item Optimizer: Adam (lr=0.001), ReduceLROnPlateau (factor=0.5, patience=10)
\end{itemize}
V1 (99 epochs, no regularization, val\_MAE=1.7516) and V3 (optimizer reset from warm-start, val\_MAE=1.0143 at epoch 5) were discarded. V2 achieves val\_MAE=0.8160 (training checkpoint) and test\_MAE=0.7985 after full evaluation.

%%====================================================================
%% FIGURE PLACEHOLDER D: LSTM V2 Training Learning Curve
%% Sumber gambar: notebooks/fase_8_modeling/08e_rul_lstm.ipynb
%% Cari cell yang memplot history['loss'] dan history['val_loss'] per epoch.
%% Export gambar sebagai: figures/lstm_learning_curve.png
%% Cara di Overleaf: Upload ke folder "figures/", uncomment blok di bawah.
%%====================================================================
% \begin{figure}[htbp]
%   \centering
%   \includegraphics[width=\columnwidth]{figures/lstm_learning_curve.png}
%   \caption{LSTM V2 training and validation MAE per epoch over 200
%            epochs. Best checkpoint at epoch~184
%            (val\_MAE~=~0.8160). ReduceLROnPlateau events are marked
%            where validation MAE plateaus.}
%   \label{fig:lstm_learning_curve}
% \end{figure}
%%====================================================================

\textbf{\textit{GRU}:} A two-layer GRU network with GRU(48) + GRU(24), Dropout(0.3), L2(0.001), BatchNormalization, Dense(16)$\rightarrow$Dense(1), totaling $\sim$27,000 parameters. Early stopping triggered at epoch 151 (of 300). Despite computational efficiency, GRU overspecializes on CRITICAL (MAE=0.014 days) while performing poorly on WARNING (MAE=1.91 days), indicating insufficient representational capacity for WARNING zone temporal complexity.

%=============================================================
\section{EXPERIMENTAL RESULTS}
%=============================================================

\subsection{Model 1: Classification Results}

Table~\ref{tab:clf_results_a} and Table~\ref{tab:clf_results_b} present the comprehensive comparison of all three classifiers, split into accuracy metrics and efficiency metrics respectively. The primary decision metric is \textbf{WARNING F1 on the validation set} (not overall F1 on the test set), as false alarms directly translate to unnecessary maintenance interventions and operational cost in industrial deployments. Fatal errors (CRITICAL$\rightarrow$HEALTHY misclassification) were zero across all models and all splits, confirming that no safety-critical misclassification occurred.

% ------------------------------------------------------------------
% TABLE Xa: Accuracy and Per-Class F1 Scores
% ------------------------------------------------------------------
\begin{table}[htbp]
\caption{TABLE Xa: Classifier Accuracy Metrics (Val and Test Sets)}
\centering
\begin{tabular}{lcccc}
\hline
\textbf{Model} & \textbf{F1$_{val}$} & \textbf{F1$_{test}$} & \textbf{Acc$_{test}$} & \textbf{Fatal Err} \\
\hline
\textit{Rand. Forest}    & 0.9292 & 0.9914 & 0.9978 & 0 \\
\textbf{\textit{XGBoost V2}} & \textbf{0.9894} & 0.9906 & 0.9974 & \textbf{0} \\
\textit{LightGBM}        & 0.9845 & 0.9876 & 0.9956 & 0 \\
\hline
\end{tabular}
\label{tab:clf_results_a}
\end{table}

\begin{table}[htbp]
\caption{TABLE Xb: Per-Class F1 on Test Set}
\centering
\begin{tabular}{lccc}
\hline
\textbf{Model} & \textbf{F1$_{HLT}$} & \textbf{F1$_{WARN}$} & \textbf{F1$_{CRIT}$} \\
\hline
\textit{Rand. Forest}        & 0.9998 & 0.9856 & 0.9889 \\
\textbf{\textit{XGBoost V2}} & 0.9996 & 0.9856 & \textbf{0.9866} \\
\textit{LightGBM}            & 0.9980 & 0.9781 & 0.9865 \\
\hline
\multicolumn{4}{l}{\textbf{Selected:} XGBoost V2 (WARN F1$_{val}$=\textbf{0.9818})} \\
\hline
\end{tabular}
\label{tab:clf_results_b}
\end{table}

\begin{table}[htbp]
\caption{TABLE Xc: Classifier Efficiency Metrics}
\centering
\begin{tabular}{lcc}
\hline
\textbf{Model} & \textbf{Inference (ms)} & \textbf{Size (MB)} \\
\hline
\textit{Rand. Forest}        & 75.25          & 3.15 \\
\textbf{\textit{XGBoost V2}} & \textbf{12.76} & 1.68 \\
\textit{LightGBM}            & 3.26           & \textbf{0.13} \\
\hline
\end{tabular}
\label{tab:clf_results_c}
\end{table}

%%====================================================================
%% FIGURE PLACEHOLDER A: XGBoost V2 Learning Curve (mlogloss vs. iteration)
%% Sumber gambar: ekspor plot dari notebook
%%   notebooks/fase_8_modeling/08b_clf_xgboost.ipynb
%% Cari cell yang memplot eval_results / learning curve.
%% Export gambar sebagai: figures/xgb_learning_curve.png
%% Cara menambah di Overleaf: Upload file ke folder "figures/",
%%   lalu uncomment blok \begin{figure} di bawah ini.
%%====================================================================
% \begin{figure}[htbp]
%   \centering
%   \includegraphics[width=\columnwidth]{figures/xgb_learning_curve.png}
%   \caption{XGBoost V2 training and validation mlogloss vs. boosting
%            iteration. Convergence at iteration 499 with val\_mlogloss
%            = 0.121, versus premature convergence at iteration 40
%            (V1, mlogloss = 0.317).}
%   \label{fig:xgb_learning_curve}
% \end{figure}
%%====================================================================

Table~\ref{tab:threshold} shows the effect of WARNING probability threshold tuning on XGBoost V2. The default threshold of 0.50 produces 954 false alarms (HEALTHY samples misclassified as WARNING), which is operationally unacceptable. Raising the threshold to 0.60 eliminates 953 false alarms while maintaining high WARNING recall, yielding the optimal precision-recall balance.

\begin{table}[htbp]
\caption{TABLE XI: XGBoost V2 WARNING Threshold Tuning (Validation Set)}
\centering
\begin{tabular}{ccccc}
\hline
\textbf{Threshold} & \textbf{WARN Prec.} & \textbf{WARN Rec.} & \textbf{F1$_{WARN}$} & \textbf{False Alarms} \\
\hline
0.50 (default) & 0.087 & 0.975 & 0.159 & 954 \\
\textbf{0.60 (optimal)} & \textbf{0.973} & \textbf{0.990} & \textbf{0.9818} & \textbf{1} \\
\hline
\end{tabular}
\label{tab:threshold}
\end{table}

\subsection{Model 2: RUL Predictor Results}

All RUL models are evaluated exclusively on WARNING+CRITICAL samples.
Table~\ref{tab:rul_results} presents the comparative results with metrics
as rows and model algorithms as columns for easy cross-model comparison.
The primary evaluation metrics are \textbf{MAE} and \textbf{Error$\leq$N~days}
(percentage of predictions within $N$ days of ground truth), as these are
operationally interpretable and directly actionable by maintenance planners.
R$^2$ and RMSE are reported for completeness but are deprioritized due to
their sensitivity to outlier RUL values (see Section~\ref{sec:discussion}).

\begin{table}[htbp]
\caption{TABLE XII: RUL Predictor Comparison---Metrics as Rows
         (WARNING+CRITICAL Only, $n_{test}$=433)}
\centering
\begin{tabular}{lp{1.5cm}p{1.5cm}p{1.5cm}}
\hline
\textbf{Metric} & \textbf{XGBoost Reg} & \textbf{LSTM V2}$^{\star}$ & \textbf{GRU} \\
\hline
MAE Val (days)   & 1.7189 & \textbf{1.6461} & 1.8618 \\
MAE Test (days)  & 1.1020 & \textbf{0.7985} & 0.9515 \\
RMSE             & 5.6002 & 6.3803          & 7.3011 \\
R$^{2}$          & 0.3961 & 0.2594          & 0.0303 \\
E$\leq$1 day     & 93.53\% & \textbf{98.04\%} & 97.80\% \\
E$\leq$3 days    & 94.92\% & \textbf{98.04\%} & 97.80\% \\
Inference        & $\sim$116ms & $\sim$116ms & $\sim$90ms \\
\hline
\multicolumn{4}{l}{$^{\star}$ \textbf{Selected model.} Best MAE and Error$\leq$1d.}\\
\hline
\end{tabular}
\label{tab:rul_results}
\end{table}

Table~\ref{tab:lstm_zone} reveals a critical asymmetry in LSTM V2 performance across prediction zones. WARNING zone predictions (208 samples) achieve near-perfect accuracy (MAE=0.10 days), reflecting the regularity of early-degradation sensor signatures. CRITICAL zone predictions (225 samples) exhibit higher error (MAE=2.03 days) driven by outlier RUL values at the tail of each machine's CRITICAL window, where sensors approach failure asymptotically.

\begin{table}[htbp]
\caption{TABLE XIII: LSTM V2 Error Analysis by Zone (Test Set)}
\centering
\begin{tabular}{lccc}
\hline
\textbf{Zone} & \textbf{N (samples)} & \textbf{MAE (days)} & \textbf{Interpretation} \\
\hline
WARNING & 208 & \textbf{0.10} & Near-perfect early prediction \\
CRITICAL & 225 & 2.03 & Tail-end outlier inflation \\
\hline
\textbf{Overall} & \textbf{433} & \textbf{0.7985} & Business-grade accuracy \\
\hline
\end{tabular}
\label{tab:lstm_zone}
\end{table}

%%====================================================================
%% FIGURE PLACEHOLDER E: LSTM V2 Predicted vs Actual RUL Scatter Plot
%% Sumber gambar: notebooks/fase_9_evaluation/09_evaluation.ipynb
%% Cari cell yang memplot scatter (y_true vs y_pred) untuk LSTM V2
%% pada test set (WARNING+CRITICAL samples).
%% Export gambar sebagai: figures/lstm_pred_vs_actual.png
%% Cara di Overleaf: Upload ke folder "figures/", uncomment blok di bawah.
%%====================================================================
% \begin{figure}[htbp]
%   \centering
%   \includegraphics[width=\columnwidth]{figures/lstm_pred_vs_actual.png}
%   \caption{LSTM V2 predicted vs. actual RUL on the test set
%            (WARNING+CRITICAL, $n$=433). Points within the $\pm$1~day
%            error band (shaded region) account for 98.04\% of
%            predictions. Outlier CRITICAL-zone points correspond to
%            tail-end RUL values exceeding 5~days.}
%   \label{fig:lstm_scatter}
% \end{figure}
%%====================================================================

\subsection{Final Model Selection and Cascade Deployment}

The cascaded deployment pipeline is defined as follows:
\begin{itemize}
\item \textbf{Model 1 (XGBoost V2, threshold=0.60):} Receives every sensor tick in real-time ($\sim$12.76ms inference). Outputs HEALTHY, WARNING, or CRITICAL.
\item \textbf{Model 2 (LSTM V2, SEQ\_LEN=24):} Activated \textit{only} when Model 1 outputs WARNING or CRITICAL. Requires a 24-timestep input window ($\sim$116ms inference including warm-up).
\item \textbf{Cascade efficiency:} Since HEALTHY constitutes $>$87\% of operational observations, Model 2 is invoked for $<$13\% of ticks on average, reducing total system inference load by approximately 70\% compared to always-on sequential inference.
\end{itemize}

The final deployed artifacts are: \texttt{classifier\_final.pkl} (1.68 MB, XGBoost V2), \texttt{rul\_predictor\_final.keras} (610.4 KB, LSTM V2), and \texttt{preprocessing\_pipeline.pkl} (6.0 KB, \texttt{StandardScaler} + feature engineering). Production smoke tests confirmed: HEALTHY inference at 46.5ms, CRITICAL with 23-entry history at 116ms, five consecutive calls averaging 49.7ms (maximum 54.0ms), all within the 200ms service-level objective.

%=============================================================
\section{DISCUSSION}
\label{sec:discussion}
%=============================================================

\textbf{SSBS vs. Random Sampling:} Standard random undersampling of the HEALTHY class inadvertently removes rows adjacent to WARNING boundaries, degrading the temporal context captured by lag and rolling features. SSBS preserves both the prime-condition reference (Block A) and the pre-failure boundary (Block B), providing balanced and contextually complete training signal. The 24-hour safety buffer in Block B prevents boundary contamination by ensuring no partially-labeled WARNING rows enter the HEALTHY sample.

\textbf{Machine-Based Split Prevents Leakage:} Global temporal splits on run-to-failure datasets systematically concentrate WARNING and CRITICAL samples into test sets (since failures occur at timeline ends), producing training sets without degradation examples. Machine-based splits guarantee that each partition contains complete failure cycles, enabling the model to learn the full degradation trajectory from each machine rather than a truncated view. This design choice directly enabled the XGBoost classifier to learn WARNING-zone patterns from 14 training machines and generalize successfully to 6 unseen machines.

\textbf{Train-Only SMOTE Prevents Temporal Contamination:} SMOTE interpolates between real minority samples in feature space. If applied before splitting, synthetic WARNING samples may encode statistical patterns derived from test-machine sensor distributions, inflating reported metrics. By strictly applying SMOTE after splitting—within the training partition only—we ensure validation and test evaluations reflect true model performance on natural, unseen data. The 6,085 synthetic samples generated had zero influence on the 3,288 validation and 2,716 test records.

\textbf{XGBoost Threshold=0.60: Precision-Recall Tradeoff:} At the default threshold of 0.50, XGBoost's raw softmax probabilities assign borderline HEALTHY samples to WARNING with non-negligible probability, producing 954 false alarms. These represent operational costs—unnecessary maintenance dispatches. Raising the threshold to 0.60 selectively excludes low-confidence WARNING predictions, reducing false alarms to 1 while maintaining near-perfect WARNING recall. This threshold was identified via validation-set precision-recall curve analysis and must be applied consistently at inference time.

\textbf{LSTM vs. GRU: WARNING Zone Complexity:} GRU (27K parameters) achieves near-perfect CRITICAL prediction (MAE=0.014 days) by learning the steep, monotonic pre-failure sensor escalation pattern. However, it fails dramatically on WARNING (MAE=1.91 days), where degradation dynamics are more gradual, multi-modal, and sensor-cross-correlated. LSTM (47K parameters) with its separate input and forget gates demonstrates superior capacity for modeling the asymmetric WARNING-zone temporal complexity. The additional 20K parameters (LSTM vs. GRU) are well justified by the 19x improvement in WARNING MAE (0.10 vs. 1.91 days).

\textbf{R$^2$ Limitation in RUL Context:} The R$^2$ coefficient of determination measures variance explained relative to the mean of the target distribution. In run-to-failure RUL datasets, a small number of CRITICAL-zone records with high RUL values (up to 146.92 days, caused by machines whose last recorded CRITICAL window preceded no subsequent failure) create outlier targets that dramatically inflate RMSE and deflate R$^2$. These outliers do not reflect actual model prediction error for the actionable prediction horizon ($<$3 days). MAE and Error$\leq$N days are resistant to these outliers and represent business-relevant accuracy metrics. LSTM V2's lower R$^2$ (0.2594 vs. XGBoost's 0.3961) does not indicate inferior operational performance, as confirmed by its superior MAE and Error$\leq$1d metrics.

\textbf{Limitations:} This study has three primary limitations. First, the dataset is synthetically generated to simulate IoT sensor behavior; real-world sensor data would exhibit additional noise characteristics (sensor drift, transmission errors, environmental interference) that may reduce model performance. Second, all 20 machines share identical type and operational context; generalization to heterogeneous machine fleets or different industries remains to be validated. Third, the pipeline is evaluated in a batch experimental setting; online learning capabilities (incremental model updates as new failure data arrives) are not implemented and constitute an important future research direction.

%=============================================================
\section{CONCLUSION}
%=============================================================

This paper presented PRIME, a complete predictive maintenance pipeline addressing the twin challenges of extreme class imbalance and temporal data leakage in industrial IoT sensor datasets. The key contributions and findings are:

\begin{itemize}
\item \textbf{SSBS} effectively reduces HEALTHY dominance from 97.36\% to 87.09\% while preserving time-series causality through chronological block selection with a 24-hour safety buffer—outperforming random undersampling by maintaining temporal integrity of lag and rolling features.

\item \textbf{Three-class Temporal Label Engineering} with an empirically derived Sensor Confirmation Layer (P90 thresholds, minimum 2 sensors) produces clean, actionable HEALTHY/WARNING/CRITICAL labels from raw binary failure signals, with window parameters ($W_{WARNING}$=48h, $W_{CRITICAL}$=24h) validated by forensic EDA across multiple machines.

\item \textbf{XGBoost V2} with probability threshold tuning (0.60) achieves WARNING F1=0.9818 on the validation set with zero fatal misclassifications (CRITICAL$\rightarrow$HEALTHY), delivering a false alarm rate reduction of 99.9\% compared to the default threshold.

\item \textbf{LSTM V2} achieves MAE=0.7985 days and 98.04\% of predictions within one day of ground truth on the test set, outperforming XGBoost Regressor by 27.5\% in MAE and 4.5 percentage points in Error$\leq$1 day.

\item The \textbf{cascaded deployment architecture} reduces real-time inference overhead by $\sim$70\% through conditional LSTM activation, demonstrating that high prediction accuracy and production efficiency are not mutually exclusive.
\end{itemize}

Future work will focus on three directions: (1) validating the pipeline on real industrial IoT datasets with authentic sensor noise and fault diversity; (2) implementing online learning mechanisms to enable continuous model adaptation as new failure data accumulates; and (3) extending the machine-based generalization framework to heterogeneous multi-factory environments with variable machine types and operational regimes.

%=============================================================
% REFERENCES
%=============================================================

\begin{thebibliography}{00}

\bibitem{lee2014industryiot}
J. Lee, B. Bagheri, and H.-A. Kao, ``A Cyber-Physical Systems architecture for Industry 4.0-based manufacturing systems,'' \textit{Manufacturing Letters}, vol. 3, pp. 18--23, 2014.

\bibitem{zhang2019pdmsurvey}
W. Zhang, D. Yang, and H. Wang, ``Data-driven methods for predictive maintenance of industrial equipment: A survey,'' \textit{IEEE Systems Journal}, vol. 13, no. 3, pp. 2213--2227, Sep. 2019.

\bibitem{hochreiter1997lstm}
S. Hochreiter and J. Schmidhuber, ``Long short-term memory,'' \textit{Neural Computation}, vol. 9, no. 8, pp. 1735--1780, Nov. 1997.

\bibitem{zheng2017lstm_rul}
S. Zheng, K. Ristovski, A. Farahat, and C. Gupta, ``Long short-term memory network for remaining useful life estimation,'' in \textit{Proc. IEEE Int. Conf. Prognostics and Health Management (ICPHM)}, Dallas, TX, USA, 2017, pp. 88--95.

\bibitem{li2018cnnlstm}
X. Li, Q. Ding, and J.-Q. Sun, ``Remaining useful life estimation in prognostics using deep convolution neural networks,'' \textit{Reliability Engineering and System Safety}, vol. 172, pp. 1--11, Apr. 2018.

\bibitem{cho2014gru}
K. Cho, B. van Merri\"{e}nboer, C. Gulcehre, D. Bahdanau, F. Bougares, H. Schwenk, and Y. Bengio, ``Learning phrase representations using RNN encoder-decoder for statistical machine translation,'' in \textit{Proc. Conf. Empirical Methods in Natural Language Processing (EMNLP)}, Doha, Qatar, 2014, pp. 1724--1734.

\bibitem{chen2016xgboost}
T. Chen and C. Guestrin, ``XGBoost: A scalable tree boosting system,'' in \textit{Proc. 22nd ACM SIGKDD Int. Conf. Knowledge Discovery and Data Mining}, San Francisco, CA, USA, 2016, pp. 785--794.

\bibitem{breiman2001rf}
L. Breiman, ``Random forests,'' \textit{Machine Learning}, vol. 45, no. 1, pp. 5--32, Oct. 2001.

\bibitem{ke2017lightgbm}
G. Ke, Q. Meng, T. Finley, T. Wang, W. Chen, W. Ma, Q. Ye, and T.-Y. Liu, ``LightGBM: A highly efficient gradient boosting decision tree,'' in \textit{Advances in Neural Information Processing Systems (NeurIPS)}, Long Beach, CA, USA, 2017, pp. 3146--3154.

\bibitem{chawla2002smote}
N. V. Chawla, K. W. Bowyer, L. O. Hall, and W. P. Kegelmeyer, ``SMOTE: Synthetic minority over-sampling technique,'' \textit{Journal of Artificial Intelligence Research}, vol. 16, pp. 321--357, Jun. 2002.

\bibitem{li2020pdmimbalance}
Y. Li, T. Li, and J. Liu, ``Machinery fault diagnosis using deep learning with class-imbalanced data in manufacturing,'' \textit{International Journal of Advanced Manufacturing Technology}, vol. 109, pp. 2447--2462, 2020.

\bibitem{blagus2013smote_bias}
R. Blagus and L. Lusa, ``SMOTE for high-dimensional class-imbalanced data,'' \textit{BMC Bioinformatics}, vol. 14, p. 106, Mar. 2013.

\bibitem{bergmeir2018tscv}
C. Bergmeir, R. J. Hyndman, and B. Koo, ``A note on the validity of cross-validation for evaluating autoregressive time series prediction,'' \textit{Computational Statistics and Data Analysis}, vol. 120, pp. 70--83, Apr. 2018.

\bibitem{kaufman2012leakage}
S. Kaufman, S. Rosset, C. Perlich, and O. Stitelman, ``Leakage in data mining: Formulation, detection, and avoidance,'' \textit{ACM Transactions on Knowledge Discovery from Data}, vol. 6, no. 4, pp. 1--21, Dec. 2012.

\bibitem{hyndman2018forecasting}
R. J. Hyndman and G. Athanasopoulos, \textit{Forecasting: Principles and Practice}, 2nd ed. Melbourne, Australia: OTexts, 2018. [Online]. Available: \url{https://otexts.com/fpp2/}

\end{thebibliography}

\end{document}
